% Options for packages loaded elsewhere
% Options for packages loaded elsewhere
\PassOptionsToPackage{unicode}{hyperref}
\PassOptionsToPackage{hyphens}{url}
\PassOptionsToPackage{dvipsnames,svgnames,x11names}{xcolor}
%
\documentclass[
  11pt,
  a4paper,
  twoside,
  openright]{scrbook}
\usepackage{xcolor}
\usepackage[top=25mm,bottom=25mm,left=30mm,right=25mm]{geometry}
\usepackage{amsmath,amssymb}
\setcounter{secnumdepth}{5}
\usepackage{iftex}
\ifPDFTeX
  \usepackage[T1]{fontenc}
  \usepackage[utf8]{inputenc}
  \usepackage{textcomp} % provide euro and other symbols
\else % if luatex or xetex
  \usepackage{unicode-math} % this also loads fontspec
  \defaultfontfeatures{Scale=MatchLowercase}
  \defaultfontfeatures[\rmfamily]{Ligatures=TeX,Scale=1}
\fi
\usepackage{lmodern}
\ifPDFTeX\else
  % xetex/luatex font selection
\fi
% Use upquote if available, for straight quotes in verbatim environments
\IfFileExists{upquote.sty}{\usepackage{upquote}}{}
\IfFileExists{microtype.sty}{% use microtype if available
  \usepackage[]{microtype}
  \UseMicrotypeSet[protrusion]{basicmath} % disable protrusion for tt fonts
}{}
\makeatletter
\@ifundefined{KOMAClassName}{% if non-KOMA class
  \IfFileExists{parskip.sty}{%
    \usepackage{parskip}
  }{% else
    \setlength{\parindent}{0pt}
    \setlength{\parskip}{6pt plus 2pt minus 1pt}}
}{% if KOMA class
  \KOMAoptions{parskip=half}}
\makeatother
% Make \paragraph and \subparagraph free-standing
\makeatletter
\ifx\paragraph\undefined\else
  \let\oldparagraph\paragraph
  \renewcommand{\paragraph}{
    \@ifstar
      \xxxParagraphStar
      \xxxParagraphNoStar
  }
  \newcommand{\xxxParagraphStar}[1]{\oldparagraph*{#1}\mbox{}}
  \newcommand{\xxxParagraphNoStar}[1]{\oldparagraph{#1}\mbox{}}
\fi
\ifx\subparagraph\undefined\else
  \let\oldsubparagraph\subparagraph
  \renewcommand{\subparagraph}{
    \@ifstar
      \xxxSubParagraphStar
      \xxxSubParagraphNoStar
  }
  \newcommand{\xxxSubParagraphStar}[1]{\oldsubparagraph*{#1}\mbox{}}
  \newcommand{\xxxSubParagraphNoStar}[1]{\oldsubparagraph{#1}\mbox{}}
\fi
\makeatother


\usepackage{longtable,booktabs,array}
\newcounter{none} % for unnumbered tables
\usepackage{calc} % for calculating minipage widths
% Correct order of tables after \paragraph or \subparagraph
\usepackage{etoolbox}
\makeatletter
\patchcmd\longtable{\par}{\if@noskipsec\mbox{}\fi\par}{}{}
\makeatother
% Allow footnotes in longtable head/foot
\IfFileExists{footnotehyper.sty}{\usepackage{footnotehyper}}{\usepackage{footnote}}
\makesavenoteenv{longtable}
\usepackage{graphicx}
\makeatletter
\newsavebox\pandoc@box
\newcommand*\pandocbounded[1]{% scales image to fit in text height/width
  \sbox\pandoc@box{#1}%
  \Gscale@div\@tempa{\textheight}{\dimexpr\ht\pandoc@box+\dp\pandoc@box\relax}%
  \Gscale@div\@tempb{\linewidth}{\wd\pandoc@box}%
  \ifdim\@tempb\p@<\@tempa\p@\let\@tempa\@tempb\fi% select the smaller of both
  \ifdim\@tempa\p@<\p@\scalebox{\@tempa}{\usebox\pandoc@box}%
  \else\usebox{\pandoc@box}%
  \fi%
}
% Set default figure placement to htbp
\def\fps@figure{htbp}
\makeatother





\setlength{\emergencystretch}{3em} % prevent overfull lines

\providecommand{\tightlist}{%
  \setlength{\itemsep}{0pt}\setlength{\parskip}{0pt}}



 


\usepackage{booktabs}
\usepackage{longtable}
\usepackage{array}
\usepackage{multirow}
\makeatletter
\@ifpackageloaded{tcolorbox}{}{\usepackage[skins,breakable]{tcolorbox}}
\@ifpackageloaded{fontawesome5}{}{\usepackage{fontawesome5}}
\definecolor{quarto-callout-color}{HTML}{909090}
\definecolor{quarto-callout-note-color}{HTML}{0758E5}
\definecolor{quarto-callout-important-color}{HTML}{CC1914}
\definecolor{quarto-callout-warning-color}{HTML}{EB9113}
\definecolor{quarto-callout-tip-color}{HTML}{00A047}
\definecolor{quarto-callout-caution-color}{HTML}{FC5300}
\definecolor{quarto-callout-color-frame}{HTML}{acacac}
\definecolor{quarto-callout-note-color-frame}{HTML}{4582ec}
\definecolor{quarto-callout-important-color-frame}{HTML}{d9534f}
\definecolor{quarto-callout-warning-color-frame}{HTML}{f0ad4e}
\definecolor{quarto-callout-tip-color-frame}{HTML}{02b875}
\definecolor{quarto-callout-caution-color-frame}{HTML}{fd7e14}
\makeatother
\makeatletter
\@ifpackageloaded{bookmark}{}{\usepackage{bookmark}}
\makeatother
\makeatletter
\@ifpackageloaded{caption}{}{\usepackage{caption}}
\AtBeginDocument{%
\ifdefined\contentsname
  \renewcommand*\contentsname{Table of contents}
\else
  \newcommand\contentsname{Table of contents}
\fi
\ifdefined\listfigurename
  \renewcommand*\listfigurename{List of Figures}
\else
  \newcommand\listfigurename{List of Figures}
\fi
\ifdefined\listtablename
  \renewcommand*\listtablename{List of Tables}
\else
  \newcommand\listtablename{List of Tables}
\fi
\ifdefined\figurename
  \renewcommand*\figurename{Figure}
\else
  \newcommand\figurename{Figure}
\fi
\ifdefined\tablename
  \renewcommand*\tablename{Table}
\else
  \newcommand\tablename{Table}
\fi
}
\@ifpackageloaded{float}{}{\usepackage{float}}
\floatstyle{ruled}
\@ifundefined{c@chapter}{\newfloat{codelisting}{h}{lop}}{\newfloat{codelisting}{h}{lop}[chapter]}
\floatname{codelisting}{Listing}
\newcommand*\listoflistings{\listof{codelisting}{List of Listings}}
\makeatother
\makeatletter
\makeatother
\makeatletter
\@ifpackageloaded{caption}{}{\usepackage{caption}}
\@ifpackageloaded{subcaption}{}{\usepackage{subcaption}}
\makeatother
\usepackage{bookmark}
\IfFileExists{xurl.sty}{\usepackage{xurl}}{} % add URL line breaks if available
\urlstyle{same}
\hypersetup{
  pdftitle={A Decision-Centred Modelling Environment for Long-Horizon Planning Under Deep Uncertainty},
  pdfauthor={Seyed Ahmad Mahmoudi Lahijani},
  colorlinks=true,
  linkcolor={0.2,0.4,0.6},
  filecolor={Maroon},
  citecolor={Blue},
  urlcolor={Blue},
  pdfcreator={LaTeX via pandoc}}


\title{A Decision-Centred Modelling Environment for Long-Horizon
Planning Under Deep Uncertainty}
\usepackage{etoolbox}
\makeatletter
\providecommand{\subtitle}[1]{% add subtitle to \maketitle
  \apptocmd{\@title}{\par {\large #1 \par}}{}{}
}
\makeatother
\subtitle{Foundations, Architecture, and Application to Industrial
Process Heat Decarbonisation}
\author{Seyed Ahmad Mahmoudi Lahijani}
\date{2025}
\begin{document}
\frontmatter
\maketitle

\renewcommand*\contentsname{Table of contents}
{
\setcounter{tocdepth}{2}
\tableofcontents
}

\mainmatter
\bookmarksetup{startatroot}

\chapter{A Decision-Centred Modelling
Environment}\label{a-decision-centred-modelling-environment}

\bookmarksetup{startatroot}

\chapter*{Welcome}\label{welcome}
\addcontentsline{toc}{chapter}{Welcome}

\markboth{Welcome}{Welcome}

This is the online companion to the manuscript:

\textbf{A Decision-Centred Modelling Environment for Long-Horizon
Planning Under Deep Uncertainty: Foundations, Architecture, and
Application to Industrial Process Heat Decarbonisation}

\emph{Seyed Ahmad Mahmoudi Lahijani}

\begin{center}\rule{0.5\linewidth}{0.5pt}\end{center}

\section*{Quick Navigation}\label{quick-navigation}
\addcontentsline{toc}{section}{Quick Navigation}

\markright{Quick Navigation}

{\def\LTcaptype{none} % do not increment counter
\begin{longtable}[]{@{}
  >{\raggedright\arraybackslash}p{(\linewidth - 2\tabcolsep) * \real{0.5000}}
  >{\raggedright\arraybackslash}p{(\linewidth - 2\tabcolsep) * \real{0.5000}}@{}}
\toprule\noalign{}
\begin{minipage}[b]{\linewidth}\raggedright
If you want to\ldots{}
\end{minipage} & \begin{minipage}[b]{\linewidth}\raggedright
Go to\ldots{}
\end{minipage} \\
\midrule\noalign{}
\endhead
\bottomrule\noalign{}
\endlastfoot
Understand the framework's big picture &
\href{front-matter/context-declaration.qmd}{Context Declaration} \\
Know how to read this document &
\href{front-matter/navigation-guide.qmd}{Navigation Guide} \\
Start the conceptual argument &
\href{manuscript/module-0-orientation.qmd}{Module 0: Orientation} \\
See what has been built & \href{manuscript/module-6-poc.qmd}{Module 6:
The Proof of Concept} \\
Ask the framework a question & \hyperref[ai-panel]{AI Discussion
Panel} \\
Contribute an extension & \href{../CONTRIBUTING.md}{Contributing} \\
\end{longtable}
}

\begin{center}\rule{0.5\linewidth}{0.5pt}\end{center}

\section*{The Core Finding}\label{the-core-finding}
\addcontentsline{toc}{section}{The Core Finding}

\markright{The Core Finding}

\begin{quote}
When the Edendale electrification pathway is evaluated against a
structured ensemble of 64 futures, \textbf{23 of those futures} produce
grid hosting capacity exceedances that generate regional infrastructure
costs invisible from within the site boundary alone.

This result --- 36 percent of the ensemble revealing a constraint that
site-level analysis cannot see --- is the gap this framework fills, made
analytically visible.
\end{quote}

\begin{center}\rule{0.5\linewidth}{0.5pt}\end{center}

\section*{The Central Commitment}\label{the-central-commitment}
\addcontentsline{toc}{section}{The Central Commitment}

\markright{The Central Commitment}

The framework is organised around a single foundational reversal:

\begin{quote}
\emph{The decision makes the boundary; the artefact makes the
connection.}
\end{quote}

This document is a living resource. Contributions, extensions, and
critiques are welcomed. See the \hyperref[community-extension]{Community
Extension} section of the Context Declaration.

\begin{tcolorbox}[enhanced jigsaw, arc=.35mm, bottomrule=.15mm, bottomtitle=1mm, breakable, colback=white, colbacktitle=quarto-callout-note-color!10!white, colframe=quarto-callout-note-color-frame, coltitle=black, left=2mm, leftrule=.75mm, opacityback=0, opacitybacktitle=0.6, rightrule=.15mm, title=\textcolor{quarto-callout-note-color}{\faInfo}\hspace{0.5em}{AI-Assisted Engagement}, titlerule=0mm, toprule=.15mm, toptitle=1mm]

Upload this document to your preferred AI assistant to ask
framework-specific questions. The Context Declaration is written as a
system prompt for AI agents. The Node Declaration at the start of every
section tells AI readers what they can skip and what they must read.

\end{tcolorbox}

\part{Front Matter}

\chapter{Context Declaration}\label{context-declaration}

For Human Readers, AI Assistants, and the Research Community

\hfill\break

\chapter{Acknowledgement of AI-Assisted
Development}\label{acknowledgement-of-ai-assisted-development}

\chapter{Preamble}\label{preamble}

How to Read This Document

\hfill\break

\chapter{Navigation Guide}\label{navigation-guide}

The Node Declaration System Explained

\hfill\break

This document is structured so that both human readers and AI agents can
navigate it efficiently, processing what is relevant to their purpose
and skipping what is not. The mechanism that makes this possible is the
\textbf{Node Declaration Table}, which appears at the start of every
section and sub-module in the manuscript. This guide explains how to
read it and how to use it.

\begin{center}\rule{0.5\linewidth}{0.5pt}\end{center}

\section{1. The Node Declaration
Table}\label{the-node-declaration-table}

Every section and sub-module opens with a table in the following format:

\begin{tcolorbox}[enhanced jigsaw, arc=.35mm, bottomrule=.15mm, bottomtitle=1mm, breakable, colback=white, colbacktitle=quarto-callout-note-color!10!white, colframe=quarto-callout-note-color-frame, coltitle=black, left=2mm, leftrule=.75mm, opacityback=0, opacitybacktitle=0.6, rightrule=.15mm, title={Node Declaration --- §X.X Section Title}, titlerule=0mm, toprule=.15mm, toptitle=1mm]

{\def\LTcaptype{none} % do not increment counter
\begin{longtable}[]{@{}
  >{\raggedright\arraybackslash}p{(\linewidth - 2\tabcolsep) * \real{0.5000}}
  >{\raggedright\arraybackslash}p{(\linewidth - 2\tabcolsep) * \real{0.5000}}@{}}
\toprule\noalign{}
\begin{minipage}[b]{\linewidth}\raggedright
Field
\end{minipage} & \begin{minipage}[b]{\linewidth}\raggedright
Content
\end{minipage} \\
\midrule\noalign{}
\endhead
\bottomrule\noalign{}
\endlastfoot
\textbf{Tier} & Section or Sub-Module \\
\textbf{Status} & Implemented / Specified / Vision \\
\textbf{Assumes} & Prior sections the reader must know \\
\textbf{Contributes} & One sentence: what this node adds to the
argument \\
\textbf{Skip condition} & Precise statement of when this node can be
skipped \\
\textbf{Passes to} & Subsequent sections that depend on this one \\
\textbf{Sub-Modules here} & Sub-modules that follow this section, or
None \\
\end{longtable}
}

\end{tcolorbox}

Each field has a specific function.

\textbf{Tier} distinguishes a main-body Section from a Sub-Module.
Sections carry the primary argument. Sub-Modules carry technical depth
that supports the argument but is not required to follow it.

\textbf{Status} declares the current state of the content this node
covers. \emph{Implemented} means the described capability exists and is
operational in the proof of concept. \emph{Specified} means the
capability has been fully designed but not yet built. \emph{Vision}
means the capability has been identified as a future direction but is
not yet specified in detail.

\textbf{Assumes} lists the prior sections a reader must have processed
for this node to make sense. If the listed sections are already known
from prior reading or professional background, the reader may enter here
without reading them first.

\textbf{Contributes} is a single sentence stating what this node adds.
It answers the question: what will I understand after reading this that
I did not understand before? It is the primary navigation aid for both
human readers and AI agents assessing whether to engage with the full
node.

\textbf{Skip condition} is a precise statement of when this node may be
passed over. It is written to be testable: a reader can assess whether
the condition applies to them and act accordingly. Skip conditions are
not invitations to skip carelessly; they are declarations of what prior
knowledge makes this node redundant.

\textbf{Passes to} lists the sections that directly depend on this node.
If you skip a section, check this field in the sections that follow it:
if your skipped section appears in their Assumes field, you may need to
return.

\textbf{Sub-Modules here} lists any sub-modules that immediately follow
this section. Sub-modules are always optional unless a downstream
section explicitly lists them in its Assumes field.

\begin{center}\rule{0.5\linewidth}{0.5pt}\end{center}

\section{2. The Skip Condition: Instructions for AI
Agents}\label{the-skip-condition-instructions-for-ai-agents}

An AI agent processing this document should apply the following rules.

If the skip condition is met, the agent may limit its processing to the
first paragraph of the section and the Node Declaration Table. The first
paragraph always states the section's core claim or function; the Node
Declaration Table provides the dependency structure. Together, these are
sufficient for the agent to maintain an accurate structural map of the
document without processing every section in full.

If the skip condition is not met, the agent must process the full
section before proceeding.

Sub-modules are always optional. No downstream section depends on a
sub-module unless that sub-module is explicitly listed in the downstream
section's Assumes field. An agent may skip all sub-modules in a first
pass and return to specific ones when a query requires the technical
depth they contain.

When answering a user query, the agent should first identify which
sections are relevant by consulting the Contributes and Passes to
fields, then process those sections fully. Sections outside this set may
be processed at summary level only.

\begin{center}\rule{0.5\linewidth}{0.5pt}\end{center}

\section{3. The Four Reading Paths}\label{the-four-reading-paths}

\textbf{Sequential argument.} Read every section in module order from
Module 0 through Module 7, consulting sub-modules where the text
references them. This path gives the fullest picture and is recommended
for researchers engaging with the methodological contribution.

\textbf{Modular reference.} Identify the section most relevant to your
question using the Table of Contents purpose annotations. Check that
section's Assumes field and read any listed prior sections you do not
already know. This path is recommended for practitioners, implementers,
and domain specialists.

\textbf{Framework overview.} Read Module 0 (Orientation) and Module 7
(Extension and Synthesis) only. Module 0 provides the full architectural
map; Module 7 provides the eleven propositions that constitute the
framework's intellectual claim. Together they give the argument's
beginning and end without the internal machinery.

\textbf{AI-assisted engagement.} Upload the full document to an AI
assistant. Use the audience-specific question catalogues in the Context
Declaration to begin. The AI assistant will use the Node Declaration
Tables to navigate the document precisely in response to your questions.
You may also ask the assistant to identify which sections are relevant
to a specific question before reading them yourself.

\begin{center}\rule{0.5\linewidth}{0.5pt}\end{center}

\section{4. The Sub-Module Structure}\label{the-sub-module-structure}

Sub-modules follow immediately after the section they deepen. They are
separated from the main section by a horizontal rule and open with their
own Node Declaration Table. Each sub-module has its own skip condition
and its own Tier field set to Sub-Module.

The purpose of this arrangement is to keep the main argument
uninterrupted. A reader following the primary argument encounters the
sub-module's Node Declaration Table, assesses its skip condition, and
either proceeds into the sub-module or returns to the main sequence. The
argument does not require the sub-module; the sub-module does not
require the reader to have read the section in any particular way before
entering it.

Sub-module identifiers follow the pattern SM-X.X-Y where X.X is the
parent section number and Y is an alphabetical letter distinguishing
multiple sub-modules within the same section.

\begin{center}\rule{0.5\linewidth}{0.5pt}\end{center}

\section{5. The Table of Contents}\label{the-table-of-contents}

Every entry in the Table of Contents carries a one-line purpose
annotation in italics. This annotation is the Contributes field stated
even more briefly. Status indicators appear after each entry: ✓ for
Implemented, ○ for Specified, and ◇ for Vision. Sub-modules are indented
one level below their parent section and marked with the prefix symbol ⌐
to distinguish them visually from main-body sections.

Example Table of Contents entries: Module 3: Architecture
\ldots\ldots\ldots\ldots\ldots\ldots\ldots\ldots\ldots\ldots\ldots\ldots\ldots\ldots\ldots\ldots..
45 §3.5 Artefact Families and Schema Design ✓ Seven canonical families
and five schema design principles \ldots{} 52 ⌐ SM-3.5-A Canonical
Artefact Schema Reference ✓ Complete field-level specifications for all
families \ldots\ldots.. 58 §3.6 Artefact Lifecycle, Validation, and
Acceptance Gates ✓ From provisional generation through validated
publication \ldots.. 65

\begin{center}\rule{0.5\linewidth}{0.5pt}\end{center}

\section{6. Worked Example}\label{worked-example}

\textbf{Hypothetical query:} ``I need to understand what fields the
SignalsPack artefact carries and how its schema is validated.''

The agent consults the Table of Contents and identifies two candidates:
§3.5 (Artefact Families and Schema Design) and SM-3.5-A (Canonical
Artefact Schema Reference).

The agent reads the Node Declaration for §3.5:

\begin{tcolorbox}[enhanced jigsaw, arc=.35mm, bottomrule=.15mm, bottomtitle=1mm, breakable, colback=white, colbacktitle=quarto-callout-note-color!10!white, colframe=quarto-callout-note-color-frame, coltitle=black, left=2mm, leftrule=.75mm, opacityback=0, opacitybacktitle=0.6, rightrule=.15mm, title={Node Declaration --- §3.5 Artefact Families and Schema Design}, titlerule=0mm, toprule=.15mm, toptitle=1mm]

{\def\LTcaptype{none} % do not increment counter
\begin{longtable}[]{@{}
  >{\raggedright\arraybackslash}p{(\linewidth - 2\tabcolsep) * \real{0.5000}}
  >{\raggedright\arraybackslash}p{(\linewidth - 2\tabcolsep) * \real{0.5000}}@{}}
\toprule\noalign{}
\begin{minipage}[b]{\linewidth}\raggedright
Field
\end{minipage} & \begin{minipage}[b]{\linewidth}\raggedright
Content
\end{minipage} \\
\midrule\noalign{}
\endhead
\bottomrule\noalign{}
\endlastfoot
\textbf{Tier} & Section \\
\textbf{Status} & ✓ Implemented \\
\textbf{Assumes} & §3.4 (The Thin-Waist Artefact Exchange) \\
\textbf{Contributes} & Defines the seven canonical artefact families and
establishes the five schema design principles governing their
structure \\
\textbf{Skip condition} & Skip if reader already knows the seven
artefact family names, their primary producers and consumers, and the
meaning of schema version stability \\
\textbf{Passes to} & §3.6, SM-3.5-A, Module 6 \\
\textbf{Sub-Modules here} & SM-3.5-A: Canonical Artefact Schema
Reference \\
\end{longtable}
}

\end{tcolorbox}

The skip condition is not met because the query requires specific
field-level knowledge (SignalsPack schema and validation), which the
skip condition does not cover. The agent processes §3.5 in full.

The agent then reads the Node Declaration for SM-3.5-A:

\begin{tcolorbox}[enhanced jigsaw, arc=.35mm, bottomrule=.15mm, bottomtitle=1mm, breakable, colback=white, colbacktitle=quarto-callout-note-color!10!white, colframe=quarto-callout-note-color-frame, coltitle=black, left=2mm, leftrule=.75mm, opacityback=0, opacitybacktitle=0.6, rightrule=.15mm, title={Node Declaration --- SM-3.5-A: Canonical Artefact Schema Reference}, titlerule=0mm, toprule=.15mm, toptitle=1mm]

{\def\LTcaptype{none} % do not increment counter
\begin{longtable}[]{@{}
  >{\raggedright\arraybackslash}p{(\linewidth - 2\tabcolsep) * \real{0.5000}}
  >{\raggedright\arraybackslash}p{(\linewidth - 2\tabcolsep) * \real{0.5000}}@{}}
\toprule\noalign{}
\begin{minipage}[b]{\linewidth}\raggedright
Field
\end{minipage} & \begin{minipage}[b]{\linewidth}\raggedright
Content
\end{minipage} \\
\midrule\noalign{}
\endhead
\bottomrule\noalign{}
\endlastfoot
\textbf{Tier} & Sub-Module \\
\textbf{Status} & ✓ Implemented \\
\textbf{Assumes} & §3.5 \\
\textbf{Contributes} & Complete field-level specifications for all seven
artefact families including column names, data types, units, and
validation rules \\
\textbf{Skip condition} & Skip if only the conceptual definition of
artefact families is needed; process in full if implementing the
framework or verifying schema compliance \\
\textbf{Passes to} & Module 6 §6.4 (implementation uses these schemas
directly) \\
\textbf{Sub-Modules here} & None \\
\end{longtable}
}

\end{tcolorbox}

The query requires the SignalsPack field-level specification and its
validation rules. The skip condition for SM-3.5-A is not met. The agent
processes SM-3.5-A in full and retrieves the SignalsPack schema fields
and their validation rules from the table within it.

The agent does not need to process §3.6, Module 4, or Module 5 to answer
this query.

\begin{center}\rule{0.5\linewidth}{0.5pt}\end{center}

\begin{tcolorbox}[enhanced jigsaw, arc=.35mm, bottomrule=.15mm, bottomtitle=1mm, breakable, colback=white, colbacktitle=quarto-callout-important-color!10!white, colframe=quarto-callout-important-color-frame, coltitle=black, left=2mm, leftrule=.75mm, opacityback=0, opacitybacktitle=0.6, rightrule=.15mm, title=\textcolor{quarto-callout-important-color}{\faExclamation}\hspace{0.5em}{Important}, titlerule=0mm, toprule=.15mm, toptitle=1mm]

This Navigation Guide is itself governed by a Node Declaration. It has
no skip condition: it should be read by all users before engaging with
the manuscript's main body.

\end{tcolorbox}

\part{Module 0: Orientation}

\chapter{Module 0: Orientation}\label{module-0-orientation-1}

The Framework at a Glance

\hfill\break

\part{Module 1: Foundations}

\chapter{Module 1: Foundations}\label{module-1-foundations-1}

The Decision Problem and Its Requirements

\hfill\break

\chapter{Sub-Module 1.1-A}\label{sub-module-1.1-a}

Philosophical Lineage of Deep Uncertainty

\hfill\break

\chapter{Sub-Module 1.4-A}\label{sub-module-1.4-a}

Decision Theory

\hfill\break

\chapter{Sub-Module 1.7-A}\label{sub-module-1.7-a}

Requirements Mapping

\hfill\break

\part{Module 2: Modelling Foundations}

\chapter{Module 2: Modelling
Foundations}\label{module-2-modelling-foundations-1}

Tools, Machinery, and the Logic of Exploration

\hfill\break

\chapter{Sub-Module 2.2-A}\label{sub-module-2.2-a}

Optimisation Roles

\hfill\break

\chapter{Sub-Module 2.3-A}\label{sub-module-2.3-a}

Surrogate Validation

\hfill\break

\chapter{Sub-Module 2.4-A}\label{sub-module-2.4-a}

Scenario Discovery

\hfill\break

\part{Module 3: Architecture}

\chapter{Module 3: Architecture}\label{module-3-architecture-1}

Modular Decomposition, Artefacts, and Governed Persistence

\hfill\break

\chapter{Sub-Module 3.5-A}\label{sub-module-3.5-a}

Artefact Schema

\hfill\break

\chapter{Sub-Module 3.7-A}\label{sub-module-3.7-a}

Backbone Implementation

\hfill\break

\chapter{Sub-Module 3.8-A}\label{sub-module-3.8-a}

AI/ML Governance

\hfill\break

\part{Module 4: Domain Translation}

\chapter{Module 4: Domain
Translation}\label{module-4-domain-translation-1}

Energy Systems, Process Heat, and Coupled Evaluation

\hfill\break

\chapter{Sub-Module 4.4-A}\label{sub-module-4.4-a}

System Regret

\hfill\break

\part{Module 5: The New Zealand Context}

\chapter{Module 5: The New Zealand
Context}\label{module-5-the-new-zealand-context-1}

Conditions, Constraints, and Empirical Grounding

\hfill\break

\chapter{Sub-Module 5.2-A}\label{sub-module-5.2-a}

GXP Capacity

\hfill\break

\chapter{Sub-Module 5.3-A}\label{sub-module-5.3-a}

Biomass Supply

\hfill\break

\chapter{Sub-Module 5.4-A}\label{sub-module-5.4-a}

ETS History

\hfill\break

\part{Module 6: The Proof of Concept}

\chapter{Module 6: The Proof of
Concept}\label{module-6-the-proof-of-concept-1}

Architecture, Implementation, and Results

\hfill\break

\chapter{Sub-Module 6.2-A}\label{sub-module-6.2-a}

PoC Pipeline

\hfill\break

\chapter{Sub-Module 6.3-B}\label{sub-module-6.3-b}

DemandPack Methods

\hfill\break

\chapter{Sub-Module 6.4-C}\label{sub-module-6.4-c}

Site Interface Contract

\hfill\break

\chapter{Sub-Module 6.5-D}\label{sub-module-6.5-d}

SignalsPack Contract

\hfill\break

\chapter{Sub-Module 6.6-E}\label{sub-module-6.6-e}

PyPSA Module Design

\hfill\break

\chapter{Sub-Module 6.8-F}\label{sub-module-6.8-f}

Validation Records

\hfill\break

\part{Module 7: Extension and Synthesis}

\chapter{Module 7: Extension and
Synthesis}\label{module-7-extension-and-synthesis-1}

Beyond Process Heat and the Closing Argument

\hfill\break

\bookmarksetup{startatroot}

\chapter{Glossary}\label{glossary}

Artefact Families, Schema Versions, and Key Terms

\hfill\break


\backmatter


\end{document}
